% ЗГЕНЕРОВАНО scripts/build_format_v2_chapter.py — не редагувати вручну.
% Джерело: books/ashtavakra-gita/poetychnyi/rozdil-2.md

\poemverse{2.1}{Я — усвідомлення без плями, \\ спокій, що вище за природу. \\ І весь цей час \\ дурила мене тільки вона — \\ омана, і більше ніщо.}
\poemverse{2.2}{Я освітлюю це тіло — \\ і тим освітлюю світ. \\ Тож світ — мій. \\ Або ж — аніщо.}
\poemverse{2.3}{Скинув світ разом із тілом — \\ і, самим лише вмінням, \\ невідомо звідки, \\ бачу те, що вище за мене.}
\poemverse{2.4}{Хвиля, піна, бульбашка — \\ усе це вода, і тільки. \\ Так і світ, що вийшов із мене: \\ не віддільний, як хвиля від моря.}
\poemverse{2.5}{Тканину дослідиш — \\ лишиться нитка. \\ Дослідиш світ — \\ лишиться Я. \\ І глибше немає нічого.}
\poemverse{2.6}{Цукор уже в тростині, \\ уже в її соку. \\ Так і світ — уже в мені, \\ просякнутий мною до дна.}
\poemverse{2.7}{Поки не знаєш Я — світ сяє. \\ Знаєш — і сяйво гасне. \\ Поки не знаєш мотузки — сяє змія. \\ Знаєш — і змія зникає.}
\poemverse{2.8}{Світло — не позичене: \\ світло і є я. \\ Коли сяє світ, \\ сяє тільки моє сяйво.}
\poemverse{2.9}{Світ, вигаданий незнанням, \\ сяє в мені — \\ срібло в перламутрі, \\ змій у мотузці, \\ вода в мареві променя.}
\poemverse{2.10}{Світ вийшов із мене — \\ у мені й розчиниться: \\ горщик у глині, \\ хвиля у воді, \\ браслет у золоті.}
\poemverse{2.11}{Поклін мені самому — \\ тому, кого не вбити, \\ хто стоїть, \\ коли гине світ \\ від Брахми до стеблини трави.}
\poemverse{2.12}{Поклін мені самому! \\ Я один. \\ Маю тіло — і не йду, не приходжу: \\ охопив усе — і стою.}
\poemverse{2.13}{Поклін мені самому — \\ тому, хто не торкається тілом, \\ а все тримає довго. \\ І немає нікого, \\ хто зрівнявся б зі мною в цьому.}
\poemverse{2.14}{Поклін мені самому! \\ У мене немає нічого — \\ або ж у мене є все, \\ до чого дістає слово й розум.}
\poemverse{2.15}{Пізнавач, пізнаване, знання — \\ трійця, якої нема. \\ Те, у чому вона сяє з незнання, \\ і є я — \\ незаплямований.}
\poemverse{2.16}{Корінь болю — двоїстість. \\ Інших ліків немає. \\ Все, що видно, — хибне. \\ Я один: смак свідомості, без плями.}
\poemverse{2.17}{Я — саме усвідомлення. \\ Незнанням я домалював собі межі. \\ Хто так думає безупинно, \\ стоїть у не-розрізненні. \\ Це моє стояння.}
\poemverse{2.18}{Світ наче в мені — \\ а в мені не вміщується. \\ Ні пут у мене, ні визволення: \\ омана вщухла, \\ бо втратила опору.}
\poemverse{2.19}{Світ разом із тілом — ніщо, \\ так установлено. \\ А я — тільки чиста свідомість. \\ Що ж іще домислювати?}
\poemverse{2.20}{Тіло, небо, пекло, \\ пута, визволення, страх — \\ усе це домалюване. \\ Що мені робити, \\ коли я — свідомість?}
\poemverse{2.21}{Навіть у натовпі, \\ де не бачу двоїстості, \\ все стає лісом. \\ То де ж мені шукати втіхи?}
\poemverse{2.22}{Я не тіло, і тіло не моє. \\ Я не живчик — я свідомість. \\ Ось де був вузол: \\ спрага жити.}
\poemverse{2.23}{Різнобарвні хвилі світів \\ раптом устали — \\ то в мені, безмежному океані, \\ здіймається вітер розуму.}
\poemverse{2.24}{Коли в безмежному океані \\ вітер розуму вщухає — \\ торговцеві-живчику лихо: \\ човен світу тоне.}
\poemverse{2.25}{У безмежному океані — диво: \\ хвилі живих істот \\ устають, бʼються, граються \\ й входять самі собою, \\ за власною природою.}
